\documentclass{article}
\title{Ch. 5 CW Pt. 2 -- Solns}
\author{Principles of Microeconomics}
\date{\today}
\usepackage{graphicx, amsmath}
\graphicspath{ {./images/} }

\begin{document}

\maketitle

\begin{enumerate}

\item

	\begin{enumerate}
	
	\item Increase -- the introduction of additional substitutes will increase the substitutability of Ford SUVs. That should increase the price elasticity of demand.
	
	\item Decrease -- the removal of substitutes will decrease the substitutability of Ford SUVs. That should decrease the price elasticity of demand.
	
	\item Increase -- if people think that SUVs are uniquely safe, that will decrease their substitutability, so the price elasticity of demand should increase.
	
	\item Increase -- a longer time horizon gives people more time to find substitutes and companies more time to introduce substitutes, such as the four-wheel-drive cargo van. That should increase the price elasticity of demand.
	
	\end{enumerate}
	
\item

	\begin{enumerate}
	
	\item
		\begin{gather*}
		\eta = \left| \frac{\% \Delta Q_D}{\% \Delta P} \right| \\
		\% \Delta Q = \frac{6 - 5}{\frac{6+5}{2}} = \frac{2}{11} \\
		\% \Delta P = \frac{4 - 5}{\frac{4+5}{2}} = \frac{-2}{9} \\
		\eta = \left| \frac{ \frac{2}{11} }{ \frac{-2}{9} } \right| = \frac{9}{11}
		\end{gather*}
		
	\item $\eta < 1$, so demand for gas is inelastic.
	
	\item Demand would be more elastic. A longer time horizon gives people more time to find substitutes or change their behavior so that they consume a different amount of gas. As a result, gas becomes more substitutable, and the price elasticity of demand increases.
	
	\end{enumerate}
	
\item

	\begin{enumerate}
	
	\item
		\begin{gather*}
		PES = \frac{\% \Delta Q_S}{\% \Delta P} \\
		1.6 = \frac{\% \Delta Q_S}{-0.18} \\
		-0.288 = \% \Delta Q_S
		\end{gather*}
		$Q_S$ decreases by 28.8\%.
		
	\item
		\begin{gather*}
		IED = \frac{\% \Delta Q_D}{\% \Delta I} \\
		0.6 = \frac{\% \Delta Q_D}{ \frac{42,000 - 38,000}{40,000} } \\
		0.6 = \frac{\% \Delta Q_D}{0.1} \\
		0.06 = \% \Delta Q_D
		\end{gather*}
		$Q_D$ increases by 6\%.
		
	\item The cross-price elasticity of demand would be negative because the two goods are complements. An increase in the price of iPhones will decrease the quantity demanded of portable chargers.
	
	\item Let $x = \% \Delta Q_D$.
		\begin{gather*}
		TR = PQ \\
		1.02TR = 1.05P*(1 + x)Q \\
		1.02TR = 1.05(1 + x) PQ \\
		1.02 = 1.05(1 + x) \\
		\frac{1.02}{1.05} - 1 = x \\
		-0.03 = x
		\\
		\eta = \left| \frac{-3\%}{5\%} \right| = 0.6
		\end{gather*}
	
	\end{enumerate}
	
\item 

\end{enumerate}

\end{document}